\documentclass[usepdftitle=false,aspectratio=169]{beamer}

\usepackage[utf8]{inputenc}
\usetheme{Singapore}
\usepackage{xcolor}
\setbeamertemplate{footline}[frame number]

\author[Fabian Bastin]{Fabian Bastin \\ \texorpdfstring{\url{fabian.bastin@umontreal.ca}}{fabian.bastin@umontreal.ca} \\ Université de Montréal -- CIRRELT -- IVADO -- Fin-ML}
\date{}
\title[KKT conditions]{KKT conditions\\Background material}

%\usepackage[francais]{babel}

\newtheorem{defn}{Definition}
\newtheorem{lem}{Lemma}
\newtheorem{thm}{Theorem}
\newtheorem{coro}{Corollary}

\def\red{\color{red}}
\def\blue{\color{blue}}

\input{../sty/macros}

\begin{document}
\frame{\titlepage}

% ------------------------------------------------------------------------------------------------------------------------------------------------------

\begin{frame}
\frametitle{Basic notions}

Consider the problem
\begin{align*}
	\min_{x \in \calX} \ & f(x) \\
	\mbox{subject to } & g_i(x) \leq 0,\ i = 1,\ldots,m, \\
	& h_j(x) = 0,\ j = 1,\ldots,r,
\end{align*}
where $\calX \subseteq \RR^n$ is the {\red ground set} on which $f$, $g$ and $h$ are defined (often $\calX = \RR^n$; it may also hold the constraints we prefer not to dualize).

The {\red feasible set} is
$$
\calF = \{ x \in \calX \,|\, g_i(x) \leq 0,\ i = 1,\ldots,m,\ h_j(x) = 0,\ j = 1,\ldots,r \}.
$$

$x^*$ is a global minimizer if $\forall x \in \calF$, $f(x^*) \leq f(x)$.

$x^*$ is a local minimizer if
$\exists \epsilon > 0$ such that $\forall x \in \calB(x^*, \epsilon) \cap \calF$, $f(x^*) \leq f(x).$

\end{frame}

\begin{frame}
	\frametitle{Lagrangian and Lagrangian dual function}
	
We define the Lagrangian as
	$$
	L(x, \lambda, \mu) = f(x) + \sum_{i = 1}^m \lambda_i g_i(x)
	+ \sum_{j = 1}^r \mu_j h_j(x),
	$$
and the dual Lagrangian function
	$$
	\calL(\lambda, \mu) = \inf_{x \in \calX} L(x, \lambda, \mu),
	$$
the infimum being taken over the ground set, and possibly equal to $-\infty$.
	
\end{frame}

\begin{frame}
\frametitle{Lagrange multipliers: equality constraints}

Consider the mathematical program
\begin{align*}
\min_{x \in \calX} \ & f(x) \\
\mbox{subject to } & h_j(x) = 0,\ j = 1,\ldots,r,
\end{align*}
where $\calX \subset \RR^n$, $f: \RR^n \rightarrow \RR$, $h_j: \RR^n \rightarrow \RR$, $j = 1,\ldots,r$.

\mbox{}

The {\red Lagrangian} of this problem is obtained by associating a Lagrange multiplier $\mu_j$ to each constraint function $h_j$:
$$
L(x, \mu) = f(x) + \sum_{j = 1}^{r} \mu_j h_j(x).
$$

We can obtain very general conditions under which $x^*$ is an optimal solution to the optimization problem, while only basic assumptions are made over $\calX$ and the functions $f$ and $h_j$, $j = 1,\ldots,r$.

\end{frame}

\begin{frame}
\frametitle{Optimality}

\begin{thm}
Assume that the Lagrangian associated to the problem
\begin{align*}
\min_{x \in \calX} \ & f(x) \\
\mbox{s.t. } & h_j(x) = 0,\ j = 1,\ldots,r,
\end{align*}
has a local minimizer $x^* \in \calX$ when the multiplier vector $\mu$ is equal to $\mu^*$.
If $h_j(x^*) = 0$, $j = 1,\ldots,r$, then $x^*$ is a local minimizer of the constrained problem.
\end{thm}

\end{frame}

\begin{frame}
\frametitle{Optimality}

\begin{proof}
Assume by contradiction that $x^*$ is not a local minimizer of the constrained problem.
Then $\forall \epsilon > 0$, $\exists\, \overline{x} \in \calX \cap \calB(x^*, \epsilon)$ such that $h_j(\overline{x}) = 0$, $j = 1,\ldots,r$, and $f(\overline{x}) < f(x^*)$.
	
Thus, $\forall \mu$,
	$$
	\sum_{j = 1}^{r} \mu_j h_j(x^*) = \sum_{j = 1}^{r} \mu_j h_j(\overline{x}) = 0,
	$$
and
$$
f(\overline{x}) + \sum_{j = 1}^{r} \mu_j h_j(\overline{x}) < f(x^*) + \sum_{j = 1}^{r} \mu_j h_j(x^*).
$$
Taking $\mu = \mu^*$, the previous inequality contradicts that $x^*$ is a local minimizer of the Lagrangian when $\mu = \mu^*$.
\end{proof}

\end{frame}

\begin{frame}
\frametitle{Lagrange multipliers: inequality constraints}

Same setting, with inequality constraints: $\calX \subset \RR^n$, $f: \RR^n \rightarrow \RR$, $g_i: \RR^n \rightarrow \RR$, $i = 1,\ldots,m$.

\mbox{}

\begin{thm}
Assume that the Lagrangian associated to the problem
\begin{align*}
\min_{x \in \calX} \ & f(x) \\
\mbox{s.t. } & g_i(x) \leq 0,\ i = 1,\ldots,m,
\end{align*}
has a local minimizer $x^* \in \calX$ when the multiplier vector $\lambda$ is equal to $\lambda^*$.
If $g_i(x^*) \leq 0$, $\lambda^*_i \geq 0$, and $\lambda^*_i g_i(x^*) = 0$, $i = 1,\ldots,m$, then $x^*$ is a local minimizer of the constrained problem.
\end{thm}

\end{frame}

\begin{frame}
\frametitle{Lagrange multipliers: inequality constraints}
	
\begin{proof}
	As previously, assume by contradiction that $x^*$ is not a local minimizer of the constrained problem.
Then $\forall \epsilon > 0$, $\exists\, \overline{x} \in \calX \cap \calB(x^*, \epsilon)$ such that $g_i(\overline{x}) \leq 0$, $i = 1,\ldots,m$, and $f(\overline{x}) < f(x^*)$.
Therefore, for $\lambda = \lambda^* \geq 0$,
$$
\sum_{i = 1}^{m} \lambda_i g_i(\overline{x}) \leq 0 \mbox{ and }
\sum_{i = 1}^{m} \lambda_i g_i(x^*) = 0.
$$
Consequently,
$$
f(\overline{x}) + \sum_{i = 1}^{m} \lambda_i g_i(\overline{x}) < f(x^*) + \sum_{i = 1}^{m} \lambda_i g_i(x^*),
$$
contradicting that $x^*$ is a local minimizer of the Lagrangian when $\lambda = \lambda^*$.
\end{proof}

\end{frame}

\begin{frame}
\frametitle{Dual problem}

The dual problem is
\begin{align*}
\max_{\lambda \in \RR^m, \mu \in \RR^r}\ & \calL(\lambda, \mu) \\
\mbox{such that } & \lambda \geq 0.
\end{align*}

Important properties:
\begin{itemize}
\item
The dual problem is always convex, i.e. $\calL$ is always concave (even if the primal problem is not convex).
\item
The primal and dual (global) optimal values, $f^*$ and $\calL^*$, always satisfy the weak duality: $f^* \geq \calL^*$.
\item {\blue Strong duality}: under some conditions (constraint qualifications), $f^* = \calL^*$.
\end{itemize}

\end{frame}

\begin{frame}
\frametitle{Saddle point}

Since
$$
\max_{\lambda \geq 0,\ \mu} L(x, \lambda, \mu) =
\begin{cases}
f(x) & \mbox{if } x \in \calF, \\
+\infty & \mbox{otherwise,}
\end{cases}
$$
the primal problem reads $\min_{x \in \calX} \max_{\lambda \geq 0,\ \mu} L(x,\lambda,\mu)$, and the dual problem
$\max_{\lambda \geq 0,\ \mu} \min_{x \in \calX} L(x,\lambda,\mu)$: weak duality is nothing but $\max \min \leq \min \max$.

\begin{defn}[Saddle point]
$(x^*, \lambda^*, \mu^*)$, with $\lambda^* \geq 0$, is a saddle point of $L$ if
$$
L(x^*, \lambda, \mu) \leq L(x^*, \lambda^*, \mu^*) \leq L(x, \lambda^*, \mu^*)
\quad \forall\, x \in \calX,\ \forall\, \lambda \geq 0,\ \forall\, \mu.
$$
\end{defn}

$(x^*, \lambda^*, \mu^*)$ is a saddle point if and only if $x^*$ is primal optimal, $(\lambda^*, \mu^*)$ is dual optimal, and the duality gap is null.

\end{frame}

\begin{frame}
\frametitle{Working with the dual function}
\small

$\calL$ is concave, being an infimum of functions that are affine in $(\lambda, \mu)$, but it is usually {\red not differentiable}.
A supergradient is however immediate.

\begin{thm}
Let $\overline{x} \in \arg\min_{x \in \calX} L(x, \lambda, \mu)$. Then, for all $(\lambda', \mu')$,
$$
\calL(\lambda', \mu') \leq \calL(\lambda, \mu) + g(\overline{x})^T(\lambda'-\lambda) + h(\overline{x})^T(\mu'-\mu).
$$
\end{thm}
\begin{proof}
$\calL(\lambda', \mu') \leq L(\overline{x}, \lambda', \mu')
= L(\overline{x}, \lambda, \mu) + g(\overline{x})^T(\lambda'-\lambda) + h(\overline{x})^T(\mu'-\mu)$,
and $L(\overline{x}, \lambda, \mu) = \calL(\lambda, \mu)$ by definition of $\overline{x}$.
\end{proof}

The constraint values at the Lagrangian minimizer thus give a supergradient, the search direction of the {\red dual ascent} methods:
$\lambda^{\nu+1} = \max \lbrace \lambda^{\nu} + \alpha^{\nu} g(\overline{x}^{\nu}), 0 \rbrace$, the maximum projecting on $\lambda \geq 0$.
As $\calL$ is non-smooth, such a direction is not necessarily an ascent one: the step $\alpha^{\nu}$ must be chosen accordingly.

\end{frame}

\begin{frame}
\frametitle{Duality gap}

Given a primal feasible solution $x$ and a dual feasible solution $(\lambda, \mu)$, the quantity $f(x) - \calL(\lambda, \mu)$ is called the duality gap between $x$ and $(\lambda, \mu)$.
Note that
$$
f(x) - f^* \leq f(x) - \calL(\lambda, \mu).
$$
Therefore, if the duality gap is equal to 0, then $x$ is primal-optimal (and similarly, $\lambda$ and $\mu$ are dual-optimal).

\mbox{}

From an algorithmic point of view, if strong duality holds, this provides a stopping criterion: if $f(x) - \calL(\lambda, \mu) \leq \epsilon$, then $f(x) - f^* \leq \epsilon$.

\end{frame}

\begin{frame}
\frametitle{Duality gap: local case}

We can also define the dual Lagrangian function restricted to the ball $\calB(x^*, \epsilon)$:
$$
\calL_{\calB(x^*, \epsilon)}(\lambda, \mu) = \min_{x \in \calB(x^*, \epsilon)} L(x,\lambda, \mu).
$$
The weak duality still holds locally:
$$
\calL^*_{\calB(x^*, \epsilon)} \leq f(x^*).
$$
Under some conditions, the strong duality also holds:
$$
\calL^*_{\calB(x^*, \epsilon)} = f(x^*).
$$

\end{frame}

\begin{frame}
\frametitle{Duality gap: local case}

Note however that
$$
\min_{x \in \mathbb{R}^n} L(x, \lambda, \mu) 
\leq \min_{x \in \calB(x^*, \epsilon)} L(x, \lambda, \mu) 
$$
so
$$
\calL^* \leq \calL^*_{\calB(x^*, \epsilon)}.
$$

\mbox{}

Therefore, if $x^*$ is a local minimizer and the strong duality locally holds,
$$
\calL^* \leq f(x^*).
$$
The inequality can be strict.

\end{frame}

\begin{frame}
	\frametitle{Karush-Kuhn-Tucker (KKT) conditions}
	
	Consider $f, g_i, h_j \in C^1$, $i = 1,\ldots,m$, $j = 1,\ldots,r$, and the problem
	\begin{align*}
		\min_{x \in \mathbb{R}^n} \ & f(x) \\
		\mbox{s.t. } & g_i(x) \leq 0,\ i = 1,\ldots,m, \\
		& h_j(x) = 0,\ j = 1,\ldots,r.
	\end{align*}
	
	Karush-Kuhn-Tucker (KKT) conditions:
	\begin{align*}
		\nabla_x L(x,\lambda,\mu) &= 0 & \mbox{(stationarity)}\\
		\lambda_i g_i(x) &= 0 & \mbox{(complementarity)} \\
		g_i(x) \leq 0, &\ h_j(x) = 0\ \forall i,j & \mbox{(primal feasibility)} \\
		\lambda_i &\geq 0\ \forall i & \mbox{(dual feasibility)}
	\end{align*}
	
\end{frame}

\begin{frame}
	\frametitle{Necessary conditions}
	
	Let $x^*$ be a minimizer of $f$ over $\calF \cap \calB(x^*,\epsilon)$, $\epsilon > 0$, and $(\lambda^*, \mu^*)$ be a dual solution if $x$ is restricted to $\calB(x^*,\epsilon)$, with a zero duality gap (the strong duality holds). Then
	\begin{align*}
		f(x^*) &= \calL_{\calB(x^*, \epsilon)}(\lambda^*, \mu^*) \\
		&= \min_{x \in \calB(x^*,\epsilon)} \left( f(x) + \sum_{i = 1}^m \lambda_i^* g_i(x) + \sum_{j = 1}^r \mu_j^* h_j(x) \right) \\
		& \leq f(x^*) + \sum_{i = 1}^m \lambda_i^* g_i(x^*) + \sum_{j = 1}^r \mu_j^* h_j(x^*) \\
		& \leq f(x^*)
	\end{align*}
	Thus, $x^*$ is a minimizer of $L(x, \lambda^*, \mu^*)$ in $\calB(x^*,\epsilon)$, and $\nabla_x L(x^*, \lambda^*, \mu^*) = 0$.
	
	\mbox{}
	
We have obtained the stationarity conditions.
	
\end{frame}

\begin{frame}
	\frametitle{Necessary conditions}
	
	The previous inequalities also imply $\sum_{i = 1}^m \lambda_i^* g_i(x^*) = 0$ as $\sum_{i = 1}^m \lambda_i^* g_i(x^*) \leq 0$, and consequently, $\lambda_i^* g_i(x^*) = 0,\ \forall i$.
	
	\mbox{}
	
	This establishes the complementarity conditions.
	
	\mbox{}
	
	If $x^*$ is a global minimizer, we can replace $\calB(x^*,\epsilon)$ by $\RR^n$.
	
	\mbox{}
	
	We can summarize our findings in the theorem below.
\begin{thm}[KKT necessary conditions]
	If $x^*$, $(\lambda^*, \mu^*)$ are primal and dual solutions with a null duality gap, then $x^*$, $(\lambda^*, \mu^*)$ satisfy the KKT conditions.
\end{thm}

\end{frame}

\begin{frame}
\frametitle{Strong duality}

\mbox{}

The strong duality assumption often plays a key role. How to ensure that it holds?
\begin{itemize}
\item
{\red Linear programming}: it always holds, as soon as one of the two programs has a finite optimal value.
\item
{\red Convex programming}: the program is convex when $\calX$ is convex, $f$ and the $g_i$ are convex, and the $h_j$ are affine.
The {\red Slater condition} is then sufficient: there exists a feasible $x$ in the relative interior of $\calX$ with $g_i(x) < 0$ for every non-affine $g_i$.
\item
{\red Nonconvex programming}: a constraint qualification is required. The most common, and the most demanding, is the linear independence constraint qualification (LICQ).
\end{itemize}

\end{frame}

\begin{frame}
\frametitle{Nonconvex programming}

\begin{thm}[Necessary conditions]
If $x^*$ is a local solution of
\begin{align*}
\min_{x \in \calX} \ & f(x) \\
\mbox{s.t. } & g_i(x) \leq 0,\ i = 1,\ldots,m \\
& h_j(x) = 0,\ j = 1,\ldots,r,
\end{align*}
where $f$, the $g_i$, $i = 1,\ldots,m$, and the $h_j$, $j = 1,\ldots,r$, are $C^1$, and a constraint qualification condition holds at $x^*$, then
$\exists (\lambda^*, \mu^*)$ such that the KKT conditions hold at $(x^*,\lambda^*, \mu^*)$.
\end{thm}

\begin{proof}
See Nocedal \& Wright, ``Numerical Optimization'', Section 12.4.
\end{proof}

\end{frame}

\begin{frame}
	\frametitle{Sufficiency of KKT conditions}
	\small
	
	Let $(x^*, \lambda^*, \mu^*)$ satisfy the KKT conditions.
	Primal feasibility and complementarity give
	$$
	L(x^*, \lambda^*, \mu^*) = f(x^*) + \sum_{i = 1}^m \lambda_i^* g_i(x^*) + \sum_{j = 1}^r \mu_j^* h_j(x^*) = f(x^*).
	$$
	
	\mbox{}
	
	{\red In the convex case}, $L(\cdot, \lambda^*, \mu^*)$ is convex, so stationarity makes $x^*$ a {\it global} minimizer of the Lagrangian over $\calX$:
	$$
	\calL(\lambda^*, \mu^*) = L(x^*, \lambda^*, \mu^*) = f(x^*).
	$$
	The duality gap is null ({\red strong duality}) and $x^*$, $(\lambda^*, \mu^*)$ are globally optimal for the primal and the dual, respectively.

	\mbox{}

	{\red In the nonconvex case}, stationarity only makes $x^*$ a critical point of $L(\cdot, \lambda^*, \mu^*)$: it may be a local minimizer, a saddle point, or even a maximizer, and $\calL(\lambda^*, \mu^*) < f(x^*)$ remains possible.
	The KKT conditions are then necessary, under a constraint qualification, but not sufficient --- not even locally, as the next slide shows.
	% donner un exemple du dernier cas
	
	%Si $x^*$ n'est pas un minimum global, la fonction dual lagrangienne n'est pas non plus minimisée globalement. Si on considère le minimum global de la fonction duale lagrangienne, le saut de dualité n'est pas nul.
	
\end{frame}

\begin{frame}
\frametitle{KKT and the local duality gap}
\small

The necessary conditions were obtained from the {\red local} strong duality. Does the converse hold: do the KKT conditions restore a null gap on a small enough ball?

\mbox{}

As $L(x^*, \lambda^*, \mu^*) = f(x^*)$, the local gap vanishes on $\calB(x^*, \epsilon)$ {\red if and only if} $x^*$ \emph{minimizes} $L(\cdot, \lambda^*, \mu^*)$ on that ball. The KKT conditions only provide {\red stationarity}, which is strictly weaker.
A constraint qualification cannot help: it is what makes the conditions {\it necessary} at a local solution, and says nothing about what they imply.

\begin{block}{Counterexample}
$\min_{x \in \RR^2} x_1^2 - x_2^2$ subject to $h(x) = x_2 = 0$.
The feasible set is the axis $x_2 = 0$, on which $f = x_1^2$: $x^* = (0,0)$ is the unique global solution, and $L(x,\mu) = x_1^2 - x_2^2 + \mu x_2$ is stationary at $x^*$ for $\mu^* = 0$.
Regularity is not at fault: $\nabla h = (0,1) \neq 0$, so {\red LICQ holds}, and on the critical cone $\lbrace d \,|\, d_2 = 0 \rbrace$, which is not trivial, $\nabla^2_{xx}L = \operatorname{diag}(2,-2)$ is positive definite --- the {\red second-order sufficient conditions hold} too.
\end{block}

\end{frame}

\begin{frame}
\frametitle{KKT and the local duality gap (cont'd)}

Yet $L(\cdot, \mu^*) = x_1^2 - x_2^2$ has a {\red saddle point} at $x^*$, so for {\it every} $\epsilon > 0$
$$
\calL_{\calB(x^*, \epsilon)}(\mu^*) = \min_{\|x\| \leq \epsilon} \left( x_1^2-x_2^2 \right) = -\epsilon^2 < 0 = f(x^*).
$$
The gap does tend to $0$ with $\epsilon$, which makes the mistake tempting, but it vanishes on no neighbourhood.
And the failure has nothing to do with the optimality of $x^*$, which is global: it comes from the behaviour of the Lagrangian itself.

\mbox{}

A sufficient condition is $\nabla^2_{xx} L \succ 0$ on the {\red whole} space, and not only on the critical cone: above, $\operatorname{diag}(2,-2)$ is positive definite on the cone but indefinite on $\RR^2$.

\end{frame}

\begin{frame}
\frametitle{Restoring a local zero gap}

The {\red augmented Lagrangian} is precisely the standard remedy:
$$
L_{\rho}(x, \lambda, \mu) = L(x, \lambda, \mu) + \frac{\rho}{2}\|h(x)\|^2,
$$
whose Hessian at $x^*$ is $\nabla^2_{xx}L + \rho\, \nabla h(x^*)\nabla h(x^*)^T$.
Under LICQ and the second-order conditions, it becomes positive definite for $\rho$ large enough, so $x^*$ locally minimizes $L_{\rho}$ and the local gap vanishes.

\mbox{}

On the previous counterexample, $L_{\rho}(x, 0) = x_1^2 + (\rho/2 - 1)x_2^2$, of Hessian $\operatorname{diag}(2, \rho-2)$: the gap is positive for $\rho \leq 2$ and null as soon as $\rho > 2$.

\mbox{}

This is the motivation for the augmented Lagrangian methods, and in stochastic programming for {\red progressive hedging}, where the penalty is placed on the nonanticipativity constraints.

\end{frame}

\begin{frame}
\frametitle{Active set}

\begin{defn}[Active set]
The active set $\calA(x)$ of the optimization problem
\begin{align*}
\min_{x \in \calX} \ & f(x) \\
\mbox{s.t. } & g_i(x) \leq 0,\ i \in \calI \\
& h_i(x) = 0,\ i \in \calE,
\end{align*}
in a feasible point $x$ is the index set of the equality constraints and the active inequality constraints at that point:
$$
\calA(x) = \calE \cup \lbrace i \in \calI \,|\, g_i(x) = 0 \rbrace
$$
\end{defn}

\end{frame}

\begin{frame}
\frametitle{LICQ}

The most popular constraint qualification is the LICQ.

\begin{defn}[LICQ]
Given a feasible point $x$ and its active set $\calA(x)$, the linear independence constraint qualification (LICQ) holds at $x$ if the gradients of the active constraints, $\lbrace \nabla_x h_i(x),\ i \in \calE \rbrace \cup \lbrace \nabla_x g_i(x),\ i \in \calA(x) \cap \calI \rbrace$, are linearly independent.
\end{defn}

\end{frame}

\end{document}